\section*{Capítulo \ref{ch:g4:clouds-rain-snow} --- Nubes, lluvia y nieve}

\begin{solution}{exo:g4:clouds-rain-snow:1}
Una \omterm{def:g4:clouds-rain-snow:cloud}{nube} está hecha de un gran número de gotitas diminutas de agua, o de
trocitos diminutos de hielo, que flotan en el aire. La \omterm{def:g4:clouds-rain-snow:cloud}{niebla} es una \omterm{def:g4:clouds-rain-snow:cloud}{nube}
que toca el suelo.
\end{solution}

\begin{solution}{exo:g4:clouds-rain-snow:2}
(a) \omterm{def:g4:clouds-rain-snow:cloud-families}{Estrato}. (b) \omterm{def:g4:clouds-rain-snow:cloud-families}{Cirros}. (c) \omterm{def:g4:clouds-rain-snow:cloud-families}{Cúmulos}. (d) \omterm{def:g4:clouds-rain-snow:cloud-families}{Cumulonimbo}.
\end{solution}

\begin{solution}{exo:g4:clouds-rain-snow:3}
(a) \omterm{def:g4:clouds-rain-snow:precipitation}{Granizo}. (b) \omterm{def:g4:clouds-rain-snow:precipitation}{Lluvia}. (c) \omterm{def:g4:clouds-rain-snow:precipitation}{Nieve}.
\end{solution}

\begin{solution}{exo:g4:clouds-rain-snow:4}
(a) Un \omterm{def:g4:clouds-rain-snow:rain-gauge}{pluviómetro}. (b) Un termómetro. (c) Una \omterm{def:g4:clouds-rain-snow:wind-vane}{veleta}.
\end{solution}

\begin{solution}{exo:g4:clouds-rain-snow:5}
Han caído 17 milímetros de \omterm{def:g4:clouds-rain-snow:precipitation}{lluvia}. $20 - 17 = 3$: harían falta 3 milímetros
más para que el agua llegara a la marca 20.
\end{solution}

\begin{solution}{exo:g4:clouds-rain-snow:6}
El día más lluvioso fue el miércoles (12 milímetros). No llovió ni el lunes
ni el viernes. El jueves cayeron 7 milímetros.
\end{solution}

\begin{solution}{exo:g4:clouds-rain-snow:7}
Los copos de \omterm{def:g4:clouds-rain-snow:precipitation}{nieve} se forman muy arriba, en la \omterm{def:g4:clouds-rain-snow:cloud}{nube} fría. Arriba, en el
pueblo de \omterm{def:g1:landscapes-around-us:landform}{montaña}, el aire está bajo cero hasta el suelo, así que llegan
como \omterm{def:g4:clouds-rain-snow:precipitation}{nieve}. Más abajo, en el \omterm{def:g1:landscapes-around-us:landform}{valle}, el aire está por encima de
\SI{0}{\celsius}: al atravesarlo, los copos se funden y llegan al suelo
como \omterm{def:g4:clouds-rain-snow:precipitation}{lluvia}.
\end{solution}

\begin{solution}{exo:g4:clouds-rain-snow:8}
El diagrama tiene siete barras de alturas 3, 0, 1, 9, 6, 0 y 1 (sin barra
el martes ni el sábado). En la semana cayeron
$3 + 0 + 1 + 9 + 6 + 0 + 1 = 20$ milímetros. El jueves trajo 9 de los
20 milímetros: $\frac{9}{20}$ de la \omterm{def:g4:clouds-rain-snow:precipitation}{lluvia} de la semana.
\end{solution}

\begin{solution}{exo:g4:clouds-rain-snow:9}
Las paredes y los árboles resguardan de la \omterm{def:g4:clouds-rain-snow:precipitation}{lluvia} el suelo que tienen
cerca, sobre todo cuando el viento empuja la \omterm{def:g4:clouds-rain-snow:precipitation}{lluvia} de lado, y de las hojas
pueden caer gotas dentro del \omterm{def:g4:clouds-rain-snow:rain-gauge}{pluviómetro}: cerca de ellos, el \omterm{def:g4:clouds-rain-snow:rain-gauge}{pluviómetro}
recogería demasiado poco, o demasiado. El termómetro tiene que medir la
temperatura del aire; a pleno sol, lo calentaría el propio sol y marcaría
de más. La caseta blanca de lamas lo protege del sol y deja pasar el aire.
\end{solution}

\begin{solution}{exo:g4:clouds-rain-snow:10}
Una \omterm{def:g4:clouds-rain-snow:wind-vane}{veleta} señala de dónde viene el viento: el viento viene del oeste. Va
hacia el este.
\end{solution}

\begin{solution}{exo:g4:clouds-rain-snow:11}
Un milímetro de \omterm{def:g4:clouds-rain-snow:precipitation}{lluvia} sobre un metro cuadrado es un litro. Sobre un metro
cuadrado, 8 milímetros dan 8 litros. Sobre 500 metros cuadrados:
$500 \times 8 = 4000$ litros. Con regaderas de 10 litros:
$4000 \div 10 = 400$ regaderas.
\end{solution}
